\subsection{Pesos exactos y aproximantes sobre la máscara corta}
En \(H_j\) la masa del punto de reinas \(p\) en cada diagonal es a lo sumo \(C\alpha\),
pues su longitud entera es menor que \(\lceil\alpha j\rceil\). Al apagar los dos factores
diagonales, el campo en las odds de cada sitio cambia en \(O(C\alpha+j^{-1})\).
En los potenciales axiales se invierte el jacobiano de márgenes en el calibre centrado.
Su complemento de Schur es \(I-p^{\mathsf T}p\) en columnas centradas. La brecha acota
allí el inverso en norma dos, y \(jp_i\le C\) acota en norma dos las secciones de su
núcleo normalizado. La identidad
\((I-T)^{-1}=I+T+T(I-T)^{-1}T\) pasa la cota a norma de supremo para \(T=p^{\mathsf T}p\).
Con \(A_{\rm ax}\) la incidencia axial y \(S_{\rm ax}=A_{\rm ax}\operatorname{diag}(p)A_{\rm ax}^t\),
el proyector \(I-A_{\rm ax}^t S_{\rm ax}^+A_{\rm ax}\operatorname{diag}(p)\) tiene por ello norma
de supremo acotada. Para \(p^{\rm T}=p(1+z)\) se restringe \(z\) a
\(A_{\rm ax}\operatorname{diag}(p)z=0\); la derivada de \(\log p^{\rm T}+\log(1-p^{\rm T})\)
proyectada en ese espacio es \(I+O(\|z\|_\infty+j^{-1})\). Restando esa gradiente
a \(z\) queda una contracción en una bola de radio \(K(C\alpha+j^{-1})\) para \(K\) de clase.
El punto obtenido es el máximo de torres por concavidad en las filas. Si una entrada
\(t_1>1/2\), para \(\xi(t)=t(1-t)/(1-2t)\) se tiene
\(\sum_{l>1}\xi(t_l)\le\xi(1-t_1)=-\xi(t_1)\) ya que \(\xi(t)/t\) es creciente
para \(t<1/2\). Cauchy--Schwarz con \(\sum_l z_l=0\) da \(\sum_l z_l^2/\xi(t_l)\ge0\),
que es la Hessiana de la acción negativa en la fila. La estricta concavidad local
identifica la rama. El punto \(p^{\rm T}\) satisface
\(p^{\rm T}_i/p_i=1+O(C\alpha+j^{-1})\). Reservando \(p/2\le p^{\rm T}\le3p/2\), se tienen
\(j\max_i p^{\rm T}_i\le2C\), brecha \(\ge\min(1,\gamma)/2\) por Rayleigh y
\(\epsilon_A:=\max_f\sum_{i\in f}p^{\rm T}_i\le2C\alpha\) en las dos familias diagonales.
Las odds \(p^{\rm T}_{rc}(1-p^{\rm T}_{rc})\) factorizan en fila y columna en el soporte. Balancear
con el campo \(\log(1-p^{\rm T}_{rc})=O(j^{-1})\) da el escalado bistocástico exacto \(D\)
de la máscara \(A\), con \(D_i/p^{\rm T}_i=1+O(j^{-1})\). La perturbación tiene sumas absolutas
\(O(j^{-1})\) por fila y columna y el mismo orden en norma de operador.

\begin{proposicion}[Menores de las caras de activación]\label{prop:menores-llegada}
Sea \(U\subset A\) con \(v=|U|\le8\log n\). Se quitan los \(v\) sitios para fijar
el exterior \(A_0=A\setminus U\). Para una cara \(Y\subseteq U\) se activan sólo sus
pares diagonales. Un subconjunto \(I\subseteq U\) compatible con la cara y con las
restricciones axiales tiene peso \(W_I=\operatorname{per}(A_{0,-I})\), donde \(-I\)
suprime las filas y columnas de \(I\). En \(H_j\), si \(q=4v+12\), los aproximantes
de orden \(q\) construidos abajo conservan los ceros prohibidos y satisfacen
\[
 |\log(\widehat W_I/W_I)|\le\eta_v=(Cq/j)^q+Ce^{-cj}.
\]
Para \(5\le v\le8\log n\) y \(j\ge\sqrt n\), el error sumado de las proyecciones
de los logaritmos de cara es \(O(j^{-8})\).
\end{proposicion}
\begin{proof}
Cada colocación determina una única intersección \(I\) con \(U\). Borrar \(v\) entradas
perturba los márgenes en \(O(v/j)\), y suprimir \(h=|I|\) filas y columnas cambia sus
medidas normalizadas en \(O(h/j)\) en variación total. El inverso de balanceo anterior
produce en cada exterior el escalado propio \(D'\), también al prolongar la ecuación de
un índice con masa cero. Tras identificar los espacios por las raíces de las medidas, la
pérdida de brecha es \(O(\sqrt{h/j}+\sqrt v/j+(v+h)/j)\): las secciones suprimidas
cuestan \(O(\sqrt{h/j})\) en norma Hilbert--Schmidt y las \(v\) entradas, \(O(\sqrt v/j)\).
Así \(m=j-h\), \(mD'_{il}\le C'\), y \(B'=D'-\mathbf1\mathbf1^{\mathsf T}/m\)
tiene norma de operador \(\le r'<1\). Hall aplicado al soporte de \(D'\) da \(W_I>0\).

Sea \(R(z)=\sum_S z^{|S|}\prod_{e\in S}B'_e\), sumado sobre emparejamientos parciales.
Para \(T\sim\operatorname{Gamma}(m+1,1)\), los momentos de \(y=m/T\) son
\(m^k(m-k)!/m!\) hasta \(k=m\); expandir el permanente da
\(m^m\operatorname{per}D'/m!=\mathbb E R(y)\).
En una descomposición singular de \(B'\), usemos gaussianas circulares \(Z_l\) de varianza
uno y \(W=\sum_l\sigma_l|Z_l|^2\). Wick representa \(R(z)\) como la media de
\(\prod_i(1+\sqrt z\,u_i)\prod_k(1+\sqrt z\,v_k)\), donde
\(u=U'\sqrt\Sigma Z,\ v=V'\sqrt\Sigma\overline Z\). Las sumas de ambos campos son cero
y sus normas cuadradas son \(W\); por media aritmética y geométrica,
\(|R(z)|\le\mathbb E(1+|z|W/m)^m\). En \(|z|r'<1\) con margen,
\(\log\mathbb E e^{|z|W}\le C\sum_l\sigma_l\le C\sqrt m\).

Para definir el truncado extraemos el factor cuadrático de la expansión de Wick.
Resultan \(d(y)=\det(I-y^2 B'^{\mathsf T}B')^{-1/2}\) y los emparejamientos de
los términos de \((1+t)e^{-t+t^2/2}\) de grados \(\ge3\), con covarianza
\(O(I+O)^{-1}\) y \(O=\left(\begin{smallmatrix}0&yB'\\yB'^{\mathsf T}&0\end{smallmatrix}\right)\).
Con \(E\) aristas y \(s\) vértices activos el exceso \(\ell=E-s\) cumple \(E\le3\ell\).
La suma de asignaciones de índices distintos se escribe por inversión en particiones de
los \(s\) vértices. Los bloques con índices repetidos llevan el factor de Möbius
\((-1)^{d-1}(d-1)!\) para tamaño \(d\). Cada propagador vale \(O(1/m)\), por difusión
y la resolvente con dos secciones exteriores. El número de emparejamientos dividido por
\(s!\) está acotado por \(C^E E^{E-s}\); los coeficientes de cada vértice y las listas
de grados cuestan \(C^E\). La suma absoluta de las inversiones, con un \(m^{-1}\) por
coincidencia, cuesta \(\prod_{i<s}(1+i/m)\le e^{Cs}\) si \(s=O(m)\).
Luego la contribución \(V_\ell(y,1/m)\) al factor restante cumple
\(|V_\ell|\le(C\ell/m)^\ell\). Definimos \(S_q=d(1+\sum_{\ell=1}^q V_\ell)\),
como función de \(y\), en un disco fijo alrededor de uno.

Fijemos un orden auxiliar \(Q_*=\lfloor\varepsilon j\rfloor\). En la variable \(w\)
de \(R(yw)\), cada arista empieza en grado uno, de modo que \(S_{Q_*}(yw)\) concuerda
hasta grado \(Q_*\). En \(|w|\le r_1>1\) con \(r_1|y|r'<1\) ambas funciones admiten
el mayorante \(Ce^{C\sqrt m}\). Cauchy deja en \(w=1\) un resto
\(Ce^{C\sqrt m}r_1^{-Q_*}=O(e^{-cj})\). Para \(y\) fuera de una ventana fija de uno,
la cota de Wick cancela \(T^m\) en la densidad gamma. El cambio \(t'=t+W\) deja el peso
\(e^W\); si \(W\le\kappa m\) la cola gamma sigue a distancia proporcional a \(m\) de su
media. En \(W>\kappa m\) se usa \(e^{-\xi\kappa m}\mathbb E e^{(1+\xi)W}\), eligiendo
\((1+\xi)r'<1\). Ambas colas cuestan \(O(e^{-cm+C\sqrt m})\).

Para \(u=1/m\) los momentos centrados \(\mu_k=\mathbb E(y-1)^k\) satisfacen
\[
 (1-ku)\mu_{k+1}=ku(2\mu_k+\mu_{k-1}),\qquad \mu_0=1,\quad\mu_1=0
\]
por integración por partes. En particular \(\mu_{2h}\le(Ch/m)^h\) hasta \(2h\le c m\),
por inducción. La suma
\(H_q=\sum_{k=0}^{2q+1}\mu_k\,\partial_y^k S_q(1,u)/k!\)
integra el Taylor hasta grado \(2q+1\); el momento par siguiente y el radio fijo
acotan el resto por \((C(q+1)/m)^{q+1}\), incluso al extender el polinomio fuera de la
ventana. La suma de excesos de \(q+1\) a \(Q_*\) tiene la misma cota, pues la razón
de mayorantes es \(\le Ce Q_*/m<1/2\). Como \(H_q=d(1)+O(1/m)\), se puede tomar el
logaritmo uniformemente. Para \(\widehat W_I\) se restituyen \(m!/m^m\) y todos los
factores exactos de escalado de este menor. En particular su normalizador usa \(m\).
La comparación positiva se conserva al sumar los átomos de una cara. Aplicar
\(\Delta_U\) cuesta \(2^v\eta_v\). Como hay a lo sumo \(j^{2v}\) soportes y
\(Cq\le j^{1/3}\) por el umbral \(n_0(b)\), el término de potencia es
\(\le2^v j^{-8-2v/3}\) por grado. El exponencial total es
\(\exp(O(\log^2 n)-cj)\). Ambos suman \(O(j^{-8})\).
\end{proof}
